% Local formal comparisons in the exact Marenbon2007 manuscript treatment.
% Not proofs of historical-system adequacy or complete medieval original admission.
\begin{annexitem}{annex:p13.early-latin.modes}
 {thm:p13-early-latin-modes}{The three literal patterns and the term-relation comparison}
\begin{proof}
The displayed \(\sim\) is classical negation, \(\&\) conjunction,
and \(\to\) material implication in this local comparison.
Assign \(p=0,q=0\). Then \(p\&q=0\), so \(\sim(p\&q)=1\),
and \(\sim p=1\), while \(q=0\). This is a countervaluation to
the first pattern.

For the second, any valuation satisfying \(p\) has \(p=1\).
If \(q=0\), then \(\sim q=1\), and \(p\&\sim q=1\), contradicting
the other premise \(\sim(p\&\sim q)\). Therefore \(q=1\).
For the third, material implication is false exactly when its antecedent
is true and its consequent false. Thus
\(\sim(p\to\sim q)\) is true exactly when \(p=1\) and
\(\sim q=0\), equivalently \(p=1,q=1\).
It is therefore equivalent to \(p\&q\), and alone already entails \(q\).
The additional premise \(p\) is redundant.
This proves that the source's literal material translation is valid,
without identifying it with the original term-incompatibility reading.
Ordinary detachment from \(p\to q,p\) is also valid: with \(p=1\),
the conditional can be true only when \(q=1\).

To make the term comparison exact, use a two-element domain \(\{a,b\}\),
\(A=\{a,b\}\), \(B=\{b\}\).
The universal \(\forall x(Ax\to\sim Bx)\) is false at \(b\);
its outer negation is true. Also \(Aa\) is true, while \(Ba\) is false.
So outer negation of the universal connection does not force \(B\)
at an arbitrarily supplied A-subject.
By contrast, assume \(\forall x\sim(Ax\&\sim Bx)\) and \(Aa\).
Instantiating the universal at \(a\), an assumption of \(\sim Ba\)
would make its prohibited conjunction true. Hence \(Ba\) holds classically.
An existential counterinstance and universal exclusion are therefore
different term relations even when their verbal descriptions sound alike.
\end{proof}
\end{annexitem}

\begin{annexitem}{annex:p13.early-latin.universals}
 {thm:p13-early-latin-universals}{All logical branches of the source objection}
\begin{proof}
Use the exact abbreviations and premises stated in the reader.
From \(M\) and \(M\to\sim N\), modus ponens gives \(\sim N\).
Temporarily assume \(E\). From \(E\to N\), obtain \(N\),
contradicting \(\sim N\). Discharge \(E\) to get \(\sim E\).
From \(E\lor T\), consider its two cases: the E-case contradicts
\(\sim E\) and therefore yields \(T\); the T-case gives \(T\) immediately.
Disjunction elimination yields \(T\) without an undischarged case assumption.

Temporarily assume \(C\). The premise \(C\to E\) yields \(E\),
again contradicting \(\sim E\). Discharge \(C\) to get \(\sim C\).
This supplies the source's intermediate conclusions through stage 9.
No premise used so far implies that a non-corresponding thought must be false.
For a countervaluation to that extra conclusion, take
\(E=0,N=0,M=1,T=1,C=0,F=0\).
The disjunction and all three implications are true, as is \(M\);
all intermediate conclusions hold, but \(F\) is false.
Thus \(F\) is not their consequence.

Adding \(\sim C\to F\) gives \(F\) by modus ponens.
Let \(Q\) abbreviate the prescription to cease inquiry.
Set \(F=1,Q=0\) while retaining the other values above.
Even the extended descriptive premises and \(F\) remain true while \(Q\)
is false. The prescription therefore needs an additional bridge, such as
\(F\to Q\), whose normative justification is not supplied by the truth table.
If that bridge is assumed, modus ponens does yield \(Q\).
This final conditional derivation does not endorse the bridge.

The proof establishes relations among the stated propositional abbreviations.
It proves neither the ontological premises nor that abstraction always
produces an accurate thought. Boethius's rejection of stage 10 is the
reported philosophical response to the opponent's added bridge.
\end{proof}
\end{annexitem}

\begin{annexitem}{annex:p13.early-latin.categorical}
 {thm:p13-early-latin-categorical}{The inclusion chain and the failed unrestricted conversion}
\begin{proof}
If \(x\in C\), then \(x\in B\) by the first inclusion, and
\(x\in A\) by the second. Thus \(C\subseteq A\).
This also shows that reversing the order in which the two premises
are written changes no consequence, provided their predication
directions are unchanged.

For failure of general conversion take domain \(\{a,b\}\),
\(A=\{a\}\), and \(B=\{a,b\}\).
Every A is B, but \(b\) is a B that is not A.
Therefore \(A\subseteq B\) does not entail \(B\subseteq A\).
If the further condition is that A and B are coextensive,
then every B is A by that very condition, so conversion is licensed.
The role of a proprium is to provide such an additional relation;
no theorem here proves that a particular biological or metaphysical
example really is a proprium.
\end{proof}
\end{annexitem}

\begin{annexitem}{annex:p13.early-latin.contraries}
 {thm:p13-early-latin-contraries}{Every minor case and the additional partition assumption}
\begin{proof}
Instantiate \(\forall x(\sim Ax\to Bx)\) at \(a\).
With \(\sim Aa\), modus ponens gives \(Ba\).
With \(\sim Ba\), assume \(\sim Aa\); the instantiated implication
would give \(Ba\), a contradiction. Hence \(\sim\sim Aa\), and
classical double-negation elimination gives \(Aa\).
These are the two licensed minor-premise cases.

For the two unlicensed cases use a singleton domain \(\{a\}\) and
\(A=B=\{a\}\). The universal conditional is true, because its negative
antecedent is false at the sole element.
Both \(Aa\) and \(Ba\) are true, and both negative conclusions are false.
Thus neither positive minor can by itself license the other's negation.

Now assume a partition: \(A\cup B=D\), \(A\cap B=\varnothing\).
If \(Aa\), disjointness excludes \(Ba\), yielding \(\sim Ba\).
If \(Ba\), it similarly yields \(\sim Aa\).
If \(\sim Aa\), exhaustiveness places \(a\) in \(B\); if \(\sim Ba\),
it places \(a\) in \(A\).
Hence all four complementary minor cases hold under the partition.
Exhaustiveness alone is insufficient for the positive-to-negative cases,
as the singleton countermodel has exhaustive but overlapping classes.
The historical immediate-contrariety condition is therefore substantive
in this declared modern extension comparison, not merely another name
for the original one-way conditional.
\end{proof}
\end{annexitem}

\begin{annexitem}{annex:p13.early-latin.topical-warrant}
 {thm:p13-early-latin-topical-warrant}{The explicit middle and the necessary warrant}
\begin{proof}
Universally instantiate the warrant at \(x=h,y=w\):
\[
 (\operatorname{Part}(h,w)\land P(w))\to P(h).
\]
The two ground premises give its conjunctive antecedent.
Modus ponens yields \(P(h)\).
No hidden premise remains in this derivation.

For the failure without the warrant use distinct elements \(h,w\),
interpret \(P\) by \(\{w\}\), and Part by
\(\{(h,h),(w,w),(h,w)\}\).
The part relation is reflexive and transitive here; it puts the
first named element below the second. Both ground premises are true,
while \(P(h)\) is false.
Thus even a well-behaved part relation does not make an arbitrary
property of the whole a property of its part.
The missing warrant is false in this countermodel.
Stating it as a premise repairs validity, not its empirical,
metaphysical or theological truth.
\end{proof}
\end{annexitem}

\begin{annexitem}{annex:p13.early-latin.modal-scope}
 {thm:p13-early-latin-modal-scope}{Two explicit countermodels and the stronger sufficient premises}
\begin{proof}
Let \(W=\{w,v\}\), with every ordered pair accessible, and evaluate
\(\Box\alpha\) as truth of \(\alpha\) at every accessible world.
This frame is reflexive.
Make \(p\) true at \(w\) and false at \(v\).
At each world the material formula \(p\to p\) is true,
so \(\Box(p\to p)\) is true at \(w\).
But \(\Box p\) is false there because \(v\) is accessible and falsifies \(p\).
Since \(p\) is true at \(w\), \(p\to\Box p\) is false at \(w\).
This is the complete scope counterexample.

For the second claim make both \(F,P\) true at \(w\) and both
false at \(v\). The implication \(F\to P\) is true at each world,
and hence its box is true at \(w\). Actual \(F\) and actual \(P\)
are true, while \(\Box P\) is false.
Thus those premises do not entail necessity of \(P\).
On any reflexive frame, \(\Box(F\to P)\) includes the implication
at the evaluation world itself; with actual \(F\), ordinary
detachment therefore gives actual \(P\).
The reflexivity hypothesis matters: with an empty accessibility
relation, every boxed implication is vacuously true even if actual
\(F\) is true and actual \(P\) false.

If additionally \(\Box F\), take any accessible world \(u\).
The first boxed premise gives \(F\to P\) at \(u\), and
the added boxed premise gives \(F\) there. Thus \(P\) is true at \(u\).
Since \(u\) was arbitrary, \(\Box P\) follows.
These are local semantic results. They do not prove that a fixed past
belief, prophecy or knowledge state legitimately supplies \(\Box F\).
That is the additional temporal and epistemic dependency in the
historical prescience discussion.
\end{proof}
\end{annexitem}

\begin{annexitem}{annex:p13.early-latin.grammaticus}
 {thm:p13-early-latin-grammaticus}{All three logical-level comparisons}
\begin{proof}
For the categorical syllogism, take an arbitrary object \(a\)
and assume \(Ga\). Instantiation of \(\forall x(Gx\to Mx)\)
and detachment give \(Ma\). Instantiation of
\(\forall x(Mx\to Sx)\) then gives \(Sa\).
Discharge the \(Ga\) assumption and universally generalize the
fresh arbitrary object. This gives \(\forall x(Gx\to Sx)\).
No premise about what a word signifies entered this derivation.

For the possibility limit use the preceding two-world reflexive
frame and a fixed domain \(\{a,b\}\).
At actual \(w\), set \(G_w=M_w=S_w=\{a\}\).
At accessible \(v\), set \(G_v=\{a,b\}\), \(M_v=\{a\}\),
and \(S_v=\{a,b\}\).
The two actual universal premises and their conclusion are true at \(w\).
At \(v\), however, \(b\) is a G that is not M.
Thus \(\Diamond\exists x(Gx\land\sim Mx)\) holds at \(w\)
while the corresponding actual existential is false and the actual
\(\forall x(Gx\to Mx)\) remains true.
Possibility of a counterinstance does not by itself refute that
actual universal. It does refute its unrestricted boxed version.

For the semantic-level coexistence use a domain with distinct
word \(t\), quality \(q\), and substance \(s\).
Interpret the quality category by \(\{q\}\), the substance category
by \(\{s\}\), the signification relation by \(\{(t,q)\}\), and
appellation by \(\{(t,s)\}\).
The word signifies a quality and appellates a substance.
The categories are disjoint, and no object belongs to both.
Consequently the two relational claims do not require a contradictory
untyped predication of a single object.
If one instead asserts substance and quality of the same \(s\)
while retaining category disjointness, that assertion is inconsistent;
changing the semantic level is therefore essential, not verbal decoration.
The model witnesses compatibility of the displayed readings, not
the historical adequacy of every detail of Anselm's semantics.
\end{proof}
\end{annexitem}
